\documentclass[12pt,a4paper]{article}
\usepackage[utf8]{vietnam}
\usepackage{amsmath,amssymb}
\usepackage{tabularx}
\usepackage{tikz}
\usepackage{xcolor}
\usepackage[top=0.6cm, bottom=0.6cm, left=0.6cm, right=0.6cm]{geometry}
\usepackage[dethi]{ex_test}

\SAOPTN{widthSA=2.8,heightSA=0.62}

\setlength{\parskip}{0pt}

% Dòng kẻ chấm mờ đều đặn chuẩn viết tay (khoảng cách 6.5mm rộng rãi vừa vặn)
\newcommand{\dongke}[1]{%
  \begingroup%
  \setlength{\parskip}{0pt}%
  \baselineskip=18.5pt%
  \vspace{1.0mm}%
  \foreach \i in {1,...,#1}{%
    \noindent{\color{gray!60}\dotfill}\par%
  }%
  \vspace{0.6mm}%
  \endgroup%
}

% Định nghĩa ô bài tập tự luận trong bảng
\newcommand{\oBaiTap}[3][6]{%
  \vspace{1.0mm}%
  \noindent\textbf{#2} #3\par%
  \dongke{#1}%
}

\begin{document}
\pagestyle{empty}
% =============================================================================
% ========================= TOÁN 6 - TUẦN 07 - BUỔI 02: MÃ ĐỀ 101 =============
% ========================= MẶT TRƯỚC (TRANG 1/2) ==============================
% =============================================================================

\noindent
\begin{tabularx}{\linewidth}{|X|c|}
\hline
\textbf{PHIẾU KIỂM TRA ĐỊNH KỲ TOÁN 6 - TUẦN 07 - BUỔI 02} & \textbf{Mã đề: 101 - Trang 1/2} \\
\hline
\end{tabularx}

\vspace{1mm}

\noindent
\begin{tabularx}{\linewidth}{|p{3.2cm}|X|}
\hline
\centering\textbf{ĐIỂM SỐ} & \centering\textbf{LỜI NHẬN XÉT CỦA THẦY CÔ} \tabularnewline
\hline
\rule{0pt}{1.8cm} & \tabularnewline
\hline
\end{tabularx}

\vspace{1.2mm}

\noindent\textbf{Câu 1 (0,25 điểm).} Tập hợp $B = \{x \in \mathbb{N} \mid 16 \le x \le 55\}$ có bao nhiêu phần tử? \shortans{40}\par\vspace{0.3mm}
\noindent\textbf{Câu 2 (0,25 điểm).} Số La Mã $\text{XIX}$ có giá trị trong hệ thập phân bằng bao nhiêu? \shortans{19}\par\vspace{0.3mm}
\noindent\textbf{Câu 3 (0,25 điểm).} Viết gọn thương $3^8 : 3^3$ dưới dạng một lũy thừa. \shortans{3^5}\par\vspace{0.3mm}
\noindent\textbf{Câu 4 (0,25 điểm).} Giá trị của biểu thức $50 - 3 \cdot 3^2$ bằng bao nhiêu? \shortans{23}\par\vspace{0.3mm}
\noindent\textbf{Câu 5 (0,25 điểm).} Số dư của phép chia tổng $(36 + 180 + 22)$ cho $9$ bằng bao nhiêu? \shortans{4}\par\vspace{0.3mm}
\noindent\textbf{Câu 6 (0,25 điểm).} Tìm chữ số $x$ để số $\overline{3x5}$ chia hết cho $9$. \shortans{1}\par\vspace{0.3mm}
\noindent\textbf{Câu 7 (0,25 điểm).} Trong các số $17; 21; 25; 27$, số nguyên tố là số nào? \shortans{17}\par\vspace{0.3mm}
\noindent\textbf{Câu 8 (0,25 điểm).} Một hình tam giác đều có bao nhiêu trục đối xứng? \shortans{3}\par\vspace{0.3mm}
\noindent\textbf{Câu 9 (0,25 điểm).} Tính chu vi mảnh đất hình vuông có cạnh $15\text{ m}$. \shortans{60\text{ m}}\par\vspace{0.3mm}
\noindent\textbf{Câu 10 (0,25 điểm).} Tính diện tích hình bình hành có đáy $12\text{ cm}$ và chiều cao $6\text{ cm}$. \shortans{72\text{ cm}^2}\par\vspace{0.3mm}
\vspace{0.4mm}

\noindent\textbf{Câu 11 (2,0 điểm). Thực hiện phép tính (hợp lý nếu có thể):}
\vspace{0.6mm}

\noindent
\begin{tabularx}{\linewidth}{|X|X|}
\hline
\oBaiTap[10]{a)}{$135 - 3^2 \cdot 4 + 2^4 \cdot 5$} & 
\oBaiTap[10]{b)}{$37 \cdot 64 + 37 \cdot 36 - 700$} \\
\hline
\oBaiTap[10]{c)}{$200 : [66 - (19 - 15)^2]$} & 
\oBaiTap[10]{d)}{$180 : \{ 75 - [22 + (16 - 14)^3] \}$} \\
\hline
\end{tabularx}

\newpage
% =============================================================================
% ========================= TOÁN 6 - TUẦN 07 - BUỔI 02: MÃ ĐỀ 101 =============
% ========================= MẶT SAU (TRANG 2/2) ================================
% =============================================================================

\noindent
\begin{tabularx}{\linewidth}{|X|c|}
\hline
\textbf{TOÁN 6 - TUẦN 07 - BUỔI 02} & \textbf{Mã đề: 101 - Trang 2/2} \\
\hline
\end{tabularx}

\vspace{1.5mm}

\noindent\textbf{Câu 12 (2,0 điểm). Tìm số tự nhiên $x$, biết:}
\vspace{0.8mm}

\noindent
\begin{tabularx}{\linewidth}{|X|X|}
\hline
\oBaiTap[5]{a)}{$3 \cdot x - 14 = 34$} & 
\oBaiTap[5]{b)}{$75 - 4 \cdot (x - 2) = 43$} \\
\hline
\oBaiTap[5]{c)}{$4 \cdot 3^{x-1} + 14 = 50$} & 
\oBaiTap[5]{d)}{Tìm số tự nhiên $x$ lớn nhất, biết $36 \ \vdots \ x$ và $54 \ \vdots \ x$.} \\
\hline
\end{tabularx}

\vspace{1.5mm}

\noindent
\begin{tabularx}{\linewidth}{|X|}
\hline
\vspace{0.5mm}
\textbf{Câu 13 (2,5 điểm).} Bác Hải có một căn phòng hình chữ nhật có chiều dài $12\text{ m}$ và chiều rộng $8\text{ m}$.\par
\textbf{a) (1,25 điểm)} Tính chu vi và diện tích của căn phòng hình chữ nhật đó.\par
\textbf{b) (1,25 điểm)} Bác Hải dùng những viên gạch men hình vuông có cạnh $40\text{ cm}$ để lát kín nền căn phòng đó. Hỏi bác Hải cần bao nhiêu viên gạch (bỏ qua mép vữa)? Biết mỗi viên gạch có giá $22\ 000$ đồng, tính tổng số tiền bác Hải phải trả để mua gạch.\par
\vspace{0.4mm}
\noindent\textit{Bài giải:}\dongke{11}
\\ \hline
\end{tabularx}

\vspace{1.5mm}

\noindent
\begin{tabularx}{\linewidth}{|X|}
\hline
\vspace{0.5mm}
\textbf{Câu 14 (1,0 điểm).} Cho tổng $A = 2 + 2^2 + 2^3 + 2^4 + \dots + 2^{20}$. Chứng minh rằng $A$ chia hết cho $3$.\par
\vspace{0.4mm}
\noindent\textit{Bài giải:}\dongke{8}
\\ \hline
\end{tabularx}

\newpage

% =============================================================================
% ========================= TOÁN 6 - TUẦN 07 - BUỔI 02: MÃ ĐỀ 102 =============
% ========================= MẶT TRƯỚC (TRANG 1/2) ==============================
% =============================================================================

\noindent
\begin{tabularx}{\linewidth}{|X|c|}
\hline
\textbf{PHIẾU KIỂM TRA ĐỊNH KỲ TOÁN 6 - TUẦN 07 - BUỔI 02} & \textbf{Mã đề: 102 - Trang 1/2} \\
\hline
\end{tabularx}

\vspace{1mm}

\noindent
\begin{tabularx}{\linewidth}{|p{3.2cm}|X|}
\hline
\centering\textbf{ĐIỂM SỐ} & \centering\textbf{LỜI NHẬN XÉT CỦA THẦY CÔ} \tabularnewline
\hline
\rule{0pt}{1.8cm} & \tabularnewline
\hline
\end{tabularx}

\vspace{1.2mm}

\noindent\textbf{Câu 1 (0,25 điểm).} Tập hợp $B = \{x \in \mathbb{N} \mid 18 \le x \le 67\}$ có bao nhiêu phần tử? \shortans{50}\par\vspace{0.3mm}
\noindent\textbf{Câu 2 (0,25 điểm).} Số La Mã $\text{XXIII}$ có giá trị trong hệ thập phân bằng bao nhiêu? \shortans{23}\par\vspace{0.3mm}
\noindent\textbf{Câu 3 (0,25 điểm).} Viết gọn thương $5^7 : 5^4$ dưới dạng một lũy thừa. \shortans{5^3}\par\vspace{0.3mm}
\noindent\textbf{Câu 4 (0,25 điểm).} Giá trị của biểu thức $70 - 4 \cdot 3^2$ bằng bao nhiêu? \shortans{34}\par\vspace{0.3mm}
\noindent\textbf{Câu 5 (0,25 điểm).} Số dư của phép chia tổng $(45 + 270 + 32)$ cho $9$ bằng bao nhiêu? \shortans{5}\par\vspace{0.3mm}
\noindent\textbf{Câu 6 (0,25 điểm).} Tìm chữ số $x$ để số $\overline{4x3}$ chia hết cho $9$. \shortans{2}\par\vspace{0.3mm}
\noindent\textbf{Câu 7 (0,25 điểm).} Trong các số $23; 33; 35; 39$, số nguyên tố là số nào? \shortans{23}\par\vspace{0.3mm}
\noindent\textbf{Câu 8 (0,25 điểm).} Một hình vuông có bao nhiêu trục đối xứng? \shortans{4}\par\vspace{0.3mm}
\noindent\textbf{Câu 9 (0,25 điểm).} Tính chu vi mảnh đất hình vuông có cạnh $18\text{ m}$. \shortans{72\text{ m}}\par\vspace{0.3mm}
\noindent\textbf{Câu 10 (0,25 điểm).} Tính diện tích hình bình hành có đáy $15\text{ cm}$ và chiều cao $8\text{ cm}$. \shortans{120\text{ cm}^2}\par\vspace{0.3mm}
\vspace{0.4mm}

\noindent\textbf{Câu 11 (2,0 điểm). Thực hiện phép tính (hợp lý nếu có thể):}
\vspace{0.6mm}

\noindent
\begin{tabularx}{\linewidth}{|X|X|}
\hline
\oBaiTap[10]{a)}{$145 - 2^4 \cdot 3 + 3^3 \cdot 2$} & 
\oBaiTap[10]{b)}{$43 \cdot 78 + 43 \cdot 22 - 300$} \\
\hline
\oBaiTap[10]{c)}{$240 : [76 - (18 - 14)^2]$} & 
\oBaiTap[10]{d)}{$220 : \{ 85 - [22 + (17 - 15)^3] \}$} \\
\hline
\end{tabularx}

\newpage
% =============================================================================
% ========================= TOÁN 6 - TUẦN 07 - BUỔI 02: MÃ ĐỀ 102 =============
% ========================= MẶT SAU (TRANG 2/2) ================================
% =============================================================================

\noindent
\begin{tabularx}{\linewidth}{|X|c|}
\hline
\textbf{TOÁN 6 - TUẦN 07 - BUỔI 02} & \textbf{Mã đề: 102 - Trang 2/2} \\
\hline
\end{tabularx}

\vspace{1.5mm}

\noindent\textbf{Câu 12 (2,0 điểm). Tìm số tự nhiên $x$, biết:}
\vspace{0.8mm}

\noindent
\begin{tabularx}{\linewidth}{|X|X|}
\hline
\oBaiTap[5]{a)}{$4 \cdot x - 17 = 39$} & 
\oBaiTap[5]{b)}{$85 - 5 \cdot (x - 3) = 45$} \\
\hline
\oBaiTap[5]{c)}{$3 \cdot 2^{x+2} + 16 = 64$} & 
\oBaiTap[5]{d)}{Tìm số tự nhiên $x$ lớn nhất, biết $40 \ \vdots \ x$ và $60 \ \vdots \ x$.} \\
\hline
\end{tabularx}

\vspace{1.5mm}

\noindent
\begin{tabularx}{\linewidth}{|X|}
\hline
\vspace{0.5mm}
\textbf{Câu 13 (2,5 điểm).} Bác Hải có một căn phòng hình chữ nhật có chiều dài $15\text{ m}$ và chiều rộng $8\text{ m}$.\par
\textbf{a) (1,25 điểm)} Tính chu vi và diện tích của căn phòng hình chữ nhật đó.\par
\textbf{b) (1,25 điểm)} Bác Hải dùng những viên gạch men hình vuông có cạnh $50\text{ cm}$ để lát kín nền căn phòng đó. Hỏi bác Hải cần bao nhiêu viên gạch (bỏ qua mép vữa)? Biết mỗi viên gạch có giá $32\ 000$ đồng, tính tổng số tiền bác Hải phải trả để mua gạch.\par
\vspace{0.4mm}
\noindent\textit{Bài giải:}\dongke{11}
\\ \hline
\end{tabularx}

\vspace{1.5mm}

\noindent
\begin{tabularx}{\linewidth}{|X|}
\hline
\vspace{0.5mm}
\textbf{Câu 14 (1,0 điểm).} Cho tổng $B = 3 + 3^2 + 3^3 + 3^4 + \dots + 3^{20}$. Chứng minh rằng $B$ chia hết cho $4$.\par
\vspace{0.4mm}
\noindent\textit{Bài giải:}\dongke{8}
\\ \hline
\end{tabularx}

\newpage

% =============================================================================
% ========================= TOÁN 6 - TUẦN 07 - BUỔI 02: MÃ ĐỀ 103 =============
% ========================= MẶT TRƯỚC (TRANG 1/2) ==============================
% =============================================================================

\noindent
\begin{tabularx}{\linewidth}{|X|c|}
\hline
\textbf{PHIẾU KIỂM TRA ĐỊNH KỲ TOÁN 6 - TUẦN 07 - BUỔI 02} & \textbf{Mã đề: 103 - Trang 1/2} \\
\hline
\end{tabularx}

\vspace{1mm}

\noindent
\begin{tabularx}{\linewidth}{|p{3.2cm}|X|}
\hline
\centering\textbf{ĐIỂM SỐ} & \centering\textbf{LỜI NHẬN XÉT CỦA THẦY CÔ} \tabularnewline
\hline
\rule{0pt}{1.8cm} & \tabularnewline
\hline
\end{tabularx}

\vspace{1.2mm}

\noindent\textbf{Câu 1 (0,25 điểm).} Tập hợp $B = \{x \in \mathbb{N} \mid 12 \le x \le 71\}$ có bao nhiêu phần tử? \shortans{60}\par\vspace{0.3mm}
\noindent\textbf{Câu 2 (0,25 điểm).} Số La Mã $\text{XXVIII}$ có giá trị trong hệ thập phân bằng bao nhiêu? \shortans{28}\par\vspace{0.3mm}
\noindent\textbf{Câu 3 (0,25 điểm).} Viết gọn thương $2^9 : 2^5$ dưới dạng một lũy thừa. \shortans{2^4}\par\vspace{0.3mm}
\noindent\textbf{Câu 4 (0,25 điểm).} Giá trị của biểu thức $80 - 2 \cdot 4^2$ bằng bao nhiêu? \shortans{48}\par\vspace{0.3mm}
\noindent\textbf{Câu 5 (0,25 điểm).} Số dư của phép chia tổng $(54 + 360 + 25)$ cho $9$ bằng bao nhiêu? \shortans{7}\par\vspace{0.3mm}
\noindent\textbf{Câu 6 (0,25 điểm).} Tìm chữ số $x$ để số $\overline{5x1}$ chia hết cho $9$. \shortans{3}\par\vspace{0.3mm}
\noindent\textbf{Câu 7 (0,25 điểm).} Trong các số $29; 39; 49; 51$, số nguyên tố là số nào? \shortans{29}\par\vspace{0.3mm}
\noindent\textbf{Câu 8 (0,25 điểm).} Một hình lục giác đều có bao nhiêu đường chéo chính? \shortans{3}\par\vspace{0.3mm}
\noindent\textbf{Câu 9 (0,25 điểm).} Tính chu vi mảnh đất hình vuông có cạnh $25\text{ m}$. \shortans{100\text{ m}}\par\vspace{0.3mm}
\noindent\textbf{Câu 10 (0,25 điểm).} Tính diện tích hình bình hành có đáy $14\text{ cm}$ và chiều cao $5\text{ cm}$. \shortans{70\text{ cm}^2}\par\vspace{0.3mm}
\vspace{0.4mm}

\noindent\textbf{Câu 11 (2,0 điểm). Thực hiện phép tính (hợp lý nếu có thể):}
\vspace{0.6mm}

\noindent
\begin{tabularx}{\linewidth}{|X|X|}
\hline
\oBaiTap[10]{a)}{$165 - 4^2 \cdot 3 + 2^5 \cdot 2$} & 
\oBaiTap[10]{b)}{$52 \cdot 83 + 52 \cdot 17 - 200$} \\
\hline
\oBaiTap[10]{c)}{$280 : [86 - (17 - 13)^2]$} & 
\oBaiTap[10]{d)}{$260 : \{ 95 - [22 + (18 - 16)^3] \}$} \\
\hline
\end{tabularx}

\newpage
% =============================================================================
% ========================= TOÁN 6 - TUẦN 07 - BUỔI 02: MÃ ĐỀ 103 =============
% ========================= MẶT SAU (TRANG 2/2) ================================
% =============================================================================

\noindent
\begin{tabularx}{\linewidth}{|X|c|}
\hline
\textbf{TOÁN 6 - TUẦN 07 - BUỔI 02} & \textbf{Mã đề: 103 - Trang 2/2} \\
\hline
\end{tabularx}

\vspace{1.5mm}

\noindent\textbf{Câu 12 (2,0 điểm). Tìm số tự nhiên $x$, biết:}
\vspace{0.8mm}

\noindent
\begin{tabularx}{\linewidth}{|X|X|}
\hline
\oBaiTap[5]{a)}{$5 \cdot x - 16 = 49$} & 
\oBaiTap[5]{b)}{$95 - 6 \cdot (x - 1) = 47$} \\
\hline
\oBaiTap[5]{c)}{$5 \cdot 2^{x+1} - 18 = 62$} & 
\oBaiTap[5]{d)}{Tìm số tự nhiên $x$ lớn nhất, biết $42 \ \vdots \ x$ và $70 \ \vdots \ x$.} \\
\hline
\end{tabularx}

\vspace{1.5mm}

\noindent
\begin{tabularx}{\linewidth}{|X|}
\hline
\vspace{0.5mm}
\textbf{Câu 13 (2,5 điểm).} Bác Hải có một căn phòng hình chữ nhật có chiều dài $16\text{ m}$ và chiều rộng $9\text{ m}$.\par
\textbf{a) (1,25 điểm)} Tính chu vi và diện tích của căn phòng hình chữ nhật đó.\par
\textbf{b) (1,25 điểm)} Bác Hải dùng những viên gạch men hình vuông có cạnh $40\text{ cm}$ để lát kín nền căn phòng đó. Hỏi bác Hải cần bao nhiêu viên gạch (bỏ qua mép vữa)? Biết mỗi viên gạch có giá $24\ 000$ đồng, tính tổng số tiền bác Hải phải trả để mua gạch.\par
\vspace{0.4mm}
\noindent\textit{Bài giải:}\dongke{11}
\\ \hline
\end{tabularx}

\vspace{1.5mm}

\noindent
\begin{tabularx}{\linewidth}{|X|}
\hline
\vspace{0.5mm}
\textbf{Câu 14 (1,0 điểm).} Cho tổng $C = 4 + 4^2 + 4^3 + 4^4 + \dots + 4^{20}$. Chứng minh rằng $C$ chia hết cho $5$.\par
\vspace{0.4mm}
\noindent\textit{Bài giải:}\dongke{8}
\\ \hline
\end{tabularx}

\newpage

% =============================================================================
% ========================= TOÁN 6 - TUẦN 07 - BUỔI 02: MÃ ĐỀ 104 =============
% ========================= MẶT TRƯỚC (TRANG 1/2) ==============================
% =============================================================================

\noindent
\begin{tabularx}{\linewidth}{|X|c|}
\hline
\textbf{PHIẾU KIỂM TRA ĐỊNH KỲ TOÁN 6 - TUẦN 07 - BUỔI 02} & \textbf{Mã đề: 104 - Trang 1/2} \\
\hline
\end{tabularx}

\vspace{1mm}

\noindent
\begin{tabularx}{\linewidth}{|p{3.2cm}|X|}
\hline
\centering\textbf{ĐIỂM SỐ} & \centering\textbf{LỜI NHẬN XÉT CỦA THẦY CÔ} \tabularnewline
\hline
\rule{0pt}{1.8cm} & \tabularnewline
\hline
\end{tabularx}

\vspace{1.2mm}

\noindent\textbf{Câu 1 (0,25 điểm).} Tập hợp $B = \{x \in \mathbb{N} \mid 15 \le x \le 84\}$ có bao nhiêu phần tử? \shortans{70}\par\vspace{0.3mm}
\noindent\textbf{Câu 2 (0,25 điểm).} Số La Mã $\text{XVI}$ có giá trị trong hệ thập phân bằng bao nhiêu? \shortans{16}\par\vspace{0.3mm}
\noindent\textbf{Câu 3 (0,25 điểm).} Viết gọn thương $7^6 : 7^2$ dưới dạng một lũy thừa. \shortans{7^4}\par\vspace{0.3mm}
\noindent\textbf{Câu 4 (0,25 điểm).} Giá trị của biểu thức $90 - 3 \cdot 4^2$ bằng bao nhiêu? \shortans{42}\par\vspace{0.3mm}
\noindent\textbf{Câu 5 (0,25 điểm).} Số dư của phép chia tổng $(63 + 180 + 42)$ cho $9$ bằng bao nhiêu? \shortans{6}\par\vspace{0.3mm}
\noindent\textbf{Câu 6 (0,25 điểm).} Tìm chữ số $x$ để số $\overline{2x3}$ chia hết cho $9$. \shortans{4}\par\vspace{0.3mm}
\noindent\textbf{Câu 7 (0,25 điểm).} Trong các số $31; 45; 55; 57$, số nguyên tố là số nào? \shortans{31}\par\vspace{0.3mm}
\noindent\textbf{Câu 8 (0,25 điểm).} Một hình thoi có bao nhiêu cạnh bằng nhau? \shortans{4}\par\vspace{0.3mm}
\noindent\textbf{Câu 9 (0,25 điểm).} Tính chu vi mảnh đất hình vuông có cạnh $16\text{ m}$. \shortans{64\text{ m}}\par\vspace{0.3mm}
\noindent\textbf{Câu 10 (0,25 điểm).} Tính diện tích hình bình hành có đáy $16\text{ cm}$ và chiều cao $7\text{ cm}$. \shortans{112\text{ cm}^2}\par\vspace{0.3mm}
\vspace{0.4mm}

\noindent\textbf{Câu 11 (2,0 điểm). Thực hiện phép tính (hợp lý nếu có thể):}
\vspace{0.6mm}

\noindent
\begin{tabularx}{\linewidth}{|X|X|}
\hline
\oBaiTap[10]{a)}{$185 - 5^2 \cdot 2 + 3^2 \cdot 6$} & 
\oBaiTap[10]{b)}{$61 \cdot 74 + 61 \cdot 26 - 100$} \\
\hline
\oBaiTap[10]{c)}{$320 : [96 - (16 - 12)^2]$} & 
\oBaiTap[10]{d)}{$300 : \{ 105 - [22 + (19 - 17)^3] \}$} \\
\hline
\end{tabularx}

\newpage
% =============================================================================
% ========================= TOÁN 6 - TUẦN 07 - BUỔI 02: MÃ ĐỀ 104 =============
% ========================= MẶT SAU (TRANG 2/2) ================================
% =============================================================================

\noindent
\begin{tabularx}{\linewidth}{|X|c|}
\hline
\textbf{TOÁN 6 - TUẦN 07 - BUỔI 02} & \textbf{Mã đề: 104 - Trang 2/2} \\
\hline
\end{tabularx}

\vspace{1.5mm}

\noindent\textbf{Câu 12 (2,0 điểm). Tìm số tự nhiên $x$, biết:}
\vspace{0.8mm}

\noindent
\begin{tabularx}{\linewidth}{|X|X|}
\hline
\oBaiTap[5]{a)}{$6 \cdot x - 19 = 53$} & 
\oBaiTap[5]{b)}{$90 - 7 \cdot (x - 2) = 41$} \\
\hline
\oBaiTap[5]{c)}{$2 \cdot 5^{x-1} + 10 = 60$} & 
\oBaiTap[5]{d)}{Tìm số tự nhiên $x$ lớn nhất, biết $48 \ \vdots \ x$ và $72 \ \vdots \ x$.} \\
\hline
\end{tabularx}

\vspace{1.5mm}

\noindent
\begin{tabularx}{\linewidth}{|X|}
\hline
\vspace{0.5mm}
\textbf{Câu 13 (2,5 điểm).} Bác Hải có một căn phòng hình chữ nhật có chiều dài $20\text{ m}$ và chiều rộng $9\text{ m}$.\par
\textbf{a) (1,25 điểm)} Tính chu vi và diện tích của căn phòng hình chữ nhật đó.\par
\textbf{b) (1,25 điểm)} Bác Hải dùng những viên gạch men hình vuông có cạnh $50\text{ cm}$ để lát kín nền căn phòng đó. Hỏi bác Hải cần bao nhiêu viên gạch (bỏ qua mép vữa)? Biết mỗi viên gạch có giá $35\ 000$ đồng, tính tổng số tiền bác Hải phải trả để mua gạch.\par
\vspace{0.4mm}
\noindent\textit{Bài giải:}\dongke{11}
\\ \hline
\end{tabularx}

\vspace{1.5mm}

\noindent
\begin{tabularx}{\linewidth}{|X|}
\hline
\vspace{0.5mm}
\textbf{Câu 14 (1,0 điểm).} Cho tổng $D = 5 + 5^2 + 5^3 + 5^4 + \dots + 5^{20}$. Chứng minh rằng $D$ chia hết cho $6$.\par
\vspace{0.4mm}
\noindent\textit{Bài giải:}\dongke{8}
\\ \hline
\end{tabularx}

\newpage

\end{document}
